\chapter{Esercizi: Rappresentazione delle informazioni}
\label{chap:esercizi-rappresentazione-informazioni}

\begin{exercise}
\label{ex:fi-rappresentazione-notazione-posizionale}
Che cosa si intende per notazione posizionale?
\end{exercise}

In un sistema di numerazione \emph{posizionale}, il valore di ciascuna
cifra dipende dalla cifra stessa, dalla sua posizione e dalla base
adottata. Il numero rappresentato si ottiene sommando le cifre
moltiplicate per le corrispondenti potenze della base
(\cref{def:fi-valore-posizionale}).

\begin{exercise}
\label{ex:fi-rappresentazione-valore-cifra-tre}
Perché nella notazione decimale la cifra 3 ha valore diverso in 345
e in 34.5?
\end{exercise}

In \((345)_{10}\) la cifra \(3\) occupa la posizione delle centinaia
e vale \(3\cdot10^2=300\); in \((34.5)_{10}\) occupa quella delle
decine e vale \(3\cdot10^1=30\). La posizione determina dunque il
peso della cifra (\cref{sec:fi-notazione-posizionale}).

\begin{exercise}
\label{ex:fi-rappresentazione-ruolo-base}
Qual è il ruolo della base in un sistema di numerazione posizionale?
\end{exercise}

La base \(b\ge2\) determina le cifre ammesse, comprese fra \(0\) e
\(b-1\), e il peso delle posizioni, espresso mediante potenze di \(b\).
Per esempio, la cifra nella posizione immediatamente a sinistra delle
unità pesa \(10\) in base dieci e \(2\) in base due
(\cref{def:fi-valore-posizionale}).

\begin{exercise}
\label{ex:fi-rappresentazione-binario-1101}
Come si calcola il valore decimale del numero binario \((1101)_2\)?
\end{exercise}

Si sommano le potenze di due associate ai bit uguali a uno:
\[
  (1101)_2=1\cdot2^3+1\cdot2^2+0\cdot2^1+1\cdot2^0=13.
\]
Il procedimento applica la \cref{def:fi-valore-posizionale}
alla base due.

\begin{exercise}
\label{ex:fi-rappresentazione-motivo-base-binaria}
Perché i computer usano principalmente la base binaria?
\end{exercise}

I circuiti digitali permettono di distinguere in modo affidabile due
stati logici, rappresentati convenzionalmente da \(0\) e \(1\).
La codifica binaria consente quindi di memorizzare ed elaborare
i dati mediante sequenze di bit
(\cref{sec:fi-conversioni-binarie}).

\begin{exercise}
\label{ex:fi-rappresentazione-conversione-parte-intera}
Come si converte la parte intera di un numero decimale in binario?
\end{exercise}

Si effettuano divisioni successive per due, registrando il resto
\(0\) oppure \(1\) di ciascuna divisione, fino a ottenere un quoziente
nullo. I resti, letti \emph{dall'ultimo al primo}, formano la
rappresentazione binaria; lo zero si rappresenta con \(0\)
(\cref{sec:fi-conversioni-binarie}).

\begin{exercise}
\label{ex:fi-rappresentazione-conversione-parte-frazionaria}
Come si converte la parte frazionaria di un numero decimale in binario?
\end{exercise}

Si moltiplica ripetutamente la parte frazionaria per due, registrando
ogni volta la parte intera del prodotto, pari a \(0\) oppure \(1\).
Si prosegue sulla frazione residua: i bit ottenuti, letti
\emph{nell'ordine di produzione}, formano la parte frazionaria
binaria. La conversione termina quando il residuo è nullo
(\cref{sec:fi-conversioni-binarie}).

\clearpage

\begin{exercise}
\label{ex:fi-rappresentazione-frazioni-periodiche}
Perché alcuni numeri razionali hanno una rappresentazione binaria
periodica?
\end{exercise}

Per la \cref{prop:fi-frazione-binaria-finita}, una frazione ridotta
ai minimi termini ha una rappresentazione binaria finita soltanto
se il denominatore è una potenza di due. Negli altri casi la
conversione produce, dopo un eventuale tratto iniziale, una sequenza
periodica, perché un residuo non nullo finisce per ripetersi
(\cref{ex:fi-frazione-periodica}).

\begin{exercise}
\label{ex:fi-rappresentazione-configurazioni-n-bit}
Con \(n\) bit, quanti numeri naturali diversi si possono rappresentare?
\end{exercise}

Ogni bit può assumere due valori, quindi \(n\) bit generano
\(2^n\) configurazioni distinte. Nell'interpretazione senza segno
ciascuna corrisponde a un naturale diverso, compreso lo zero
(\cref{sec:fi-interi-rappresentazione}).

\begin{exercise}
\label{ex:fi-rappresentazione-intervallo-senza-segno}
Qual è l'intervallo dei numeri naturali rappresentabili con \(n\) bit?
\end{exercise}

Nella rappresentazione senza segno l'intervallo è
\[
  [0,\,2^n-1]\cap\Z.
\]
Per esempio, con otto bit si rappresentano gli interi da \(0\) a
\(255\) (\cref{sec:fi-interi-rappresentazione}).

\begin{exercise}
\label{ex:fi-rappresentazione-modulo-segno}
Come funziona la rappresentazione modulo e segno?
\end{exercise}

Il bit più significativo indica il segno: \(0\) per i valori
non negativi e \(1\) per quelli negativi. I restanti \(n-1\) bit
rappresentano il modulo del numero, che assume quindi valore
\((-1)^sM\). Il formato richiede almeno due bit
(\cref{sec:fi-interi-rappresentazione}).

\begin{exercise}
\label{ex:fi-rappresentazione-limiti-modulo-segno}
Perché la rappresentazione modulo e segno non è usata normalmente
nei computer?
\end{exercise}

Presenta due configurazioni dello zero, \(+0\) e \(-0\), e rende
l'addizione più complessa, perché richiede di distinguere i segni
degli operandi e, quando sono opposti, di confrontarne i moduli.
Il complemento a due evita la doppia codifica dello zero e semplifica
l'aritmetica intera (\cref{sec:fi-interi-rappresentazione,sec:fi-complemento-due}).

\begin{exercise}
\label{ex:fi-rappresentazione-definizione-complemento-due}
Che cos'è il complemento a due?
\end{exercise}

È una codifica degli interi con segno su \(n\) bit: i valori
non negativi conservano la rappresentazione binaria ordinaria,
mentre un valore negativo \(a\) è codificato mediante il valore
senza segno \(a+2^n\). La codifica è univoca nell'intervallo
rappresentabile (\cref{def:fi-complemento-due}).

\begin{exercise}
\label{ex:fi-rappresentazione-intervallo-complemento-due}
Qual è l'intervallo rappresentabile con \(n\) bit in complemento a due?
\end{exercise}

L'intervallo degli interi rappresentabili è
\[
  [-2^{n-1},\,2^{n-1}-1]\cap\Z.
\]
Con otto bit gli estremi sono \(-128\) e \(127\): è presente
un valore negativo in più rispetto ai positivi
(\cref{def:fi-complemento-due}).

\begin{exercise}
\label{ex:fi-rappresentazione-meno-quattordici}
Come si rappresenta \(-14\) in complemento a due su 8 bit?
\end{exercise}

Si rappresenta \(+14\) su otto bit, si invertono i bit e si aggiunge uno:
\[
  0000\,1110 \;\longrightarrow\; 1111\,0001
  \;\xrightarrow{+1}\;1111\,0010.
\]
La verifica numerica è \(2^8-14=242=(1111\,0010)_2\)
(\cref{sec:fi-complemento-due}).

\begin{exercise}
\label{ex:fi-rappresentazione-somma-complemento-due}
Come si esegue una somma tra numeri rappresentati in complemento
a due?
\end{exercise}

Si esegue la normale addizione binaria delle due configurazioni
su \(n\) bit, scartando l'eventuale riporto finale. Il risultato
corrisponde alla somma matematica modulo \(2^n\); rappresenta
correttamente la somma con segno soltanto se non si verifica overflow
(\cref{sec:fi-aritmetica-overflow}).

\begin{exercise}
\label{ex:fi-rappresentazione-overflow-somma}
Quando si verifica overflow nella somma di due numeri in
complemento a due?
\end{exercise}

Si verifica \emph{overflow} quando la somma matematica eccede
l'intervallo rappresentabile su \(n\) bit. Nell'addizione in
complemento a due, ciò accade esattamente quando gli operandi
hanno lo stesso segno e il risultato ha segno opposto
(\cref{prop:fi-overflow-somma}).

\begin{exercise}
\label{ex:fi-rappresentazione-reali-approssimati}
Perché i numeri reali devono essere rappresentati in modo
approssimato nei computer?
\end{exercise}

Un formato di lunghezza finita ammette soltanto un numero finito
di configurazioni, mentre i numeri reali sono infiniti. Perciò
non tutti i reali possono essere codificati esattamente: i valori
non rappresentabili vengono approssimati secondo il formato
adottato (\cref{sec:fi-virgola-mobile,sec:fi-precisione-arrotondamento}).

\begin{exercise}
\label{ex:fi-rappresentazione-fissa-mobile}
Che differenza c'è tra rappresentazione in virgola fissa e
rappresentazione in virgola mobile?
\end{exercise}

Nella \emph{virgola fissa} la posizione del separatore è prestabilita,
quindi i valori rappresentabili hanno un passo costante. Nella
\emph{virgola mobile} il valore dipende da un significando e da
un esponente, consentendo di coprire ordini di grandezza diversi
con una precisione relativa limitata
(\cref{sec:fi-virgola-mobile}).

\begin{exercise}
\label{ex:fi-rappresentazione-mantissa-esponente}
In una rappresentazione in virgola mobile, che cosa sono mantissa
ed esponente?
\end{exercise}

Il \emph{significando}, detto anche \emph{mantissa}, contiene le cifre
significative del numero; l'\emph{esponente} ne determina l'ordine
di grandezza, moltiplicando il significando per una potenza della
base. Il segno completa la rappresentazione, secondo la forma
\(x=(-1)^sMb^E\) (\cref{sec:fi-virgola-mobile}).

\begin{exercise}
\label{ex:fi-rappresentazione-ieee32-campi}
Quali sono i campi principali di un numero floating point a 32 bit
secondo IEEE 754?
\end{exercise}

Il formato \emph{binary32} comprende un bit di segno, otto bit
di esponente e ventitré bit per la frazione del significando.
Nei valori normali il primo bit \(1\) del significando è implicito,
per un totale di ventiquattro bit significativi
(\cref{sec:fi-ieee754,tab:fi-formati-ieee}).

\begin{exercise}
\label{ex:fi-rappresentazione-ieee32-esponenti-speciali}
Che cosa rappresentano gli esponenti speciali 00000000 e 11111111
nello standard IEEE 754?
\end{exercise}

Nel formato binary32, il campo esponente tutto a zero indica
\(\pm0\) se anche la frazione è nulla, oppure un valore subnormale
se la frazione non è nulla. Il campo tutto a uno indica \(\pm\infty\)
se la frazione è nulla, oppure NaN se non è nulla; il bit di segno
distingue i due segni degli zeri e delle infinità
(\cref{tab:fi-valori-speciali}).

\begin{exercise}
\label{ex:fi-rappresentazione-subnormali}
Che cosa sono i numeri subnormali e perché sono utili?
\end{exercise}

Sono valori non nulli con campo esponente nullo e frazione non nulla;
a differenza dei valori normali, il loro significando non possiede
il bit iniziale implicito \(1\). Permettono di rappresentare
numeri più vicini a zero del minimo normale positivo e realizzano
un \emph{underflow graduale}, a costo di una precisione relativa
minore (\cref{def:fi-valori-ieee,sec:fi-ieee754}).

\begin{exercise}
\label{ex:fi-rappresentazione-nan}
Che cosa significa NaN e in quali casi può comparire?
\end{exercise}

\emph{NaN} significa \emph{Not a Number} e indica un risultato
non numerico nell'aritmetica IEEE 754, come \(0/0\) oppure
\(+\infty-(+\infty)\). È codificato mediante un campo esponente
tutto a uno e una frazione non nulla. Il confronto di uguaglianza
con NaN è sempre falso, anche rispetto a sé stesso
(\cref{tab:fi-valori-speciali,sec:fi-ieee754}).

\begin{exercise}
\label{ex:fi-rappresentazione-arrotondamento-435}
Perché il numero 4.35 può causare errori di arrotondamento in binario?
\end{exercise}

La frazione \(0.35=7/20\) non ha una rappresentazione binaria
finita, perché il denominatore ridotto non è una potenza di due.
Il valore \(4.35\) viene quindi approssimato: in binary64 la sua
moltiplicazione per \(100\), arrotondata nel medesimo formato,
può produrre \(434.99999999999994\) anziché \(435\)
(\cref{prop:fi-frazione-binaria-finita,ex:fi-arrotondamento-prodotto}).

\begin{exercise}
\label{ex:fi-rappresentazione-esadecimale}
Che cos'è la rappresentazione esadecimale e perché è comoda
per rappresentare sequenze di bit?
\end{exercise}

La notazione \emph{esadecimale} usa la base sedici, con cifre
\(0,\ldots,9,\mathrm{A},\ldots,\mathrm{F}\). Poiché \(16=2^4\),
ogni cifra corrisponde esattamente a quattro bit: una sequenza
binaria può quindi essere scritta in modo più breve senza
perdere informazione (\cref{sec:fi-esadecimale-ottale}).

\begin{exercise}
\label{ex:fi-rappresentazione-conversione-esadecimale}
Come si converte un numero binario in esadecimale?
\end{exercise}

Si dividono i bit in gruppi di quattro, partendo dal separatore
fra parte intera e frazionaria. Ogni gruppo viene sostituito con
la cifra esadecimale corrispondente; se necessario, si aggiungono
zeri a sinistra della parte intera o a destra di quella frazionaria.
Per esempio, \((1101\,0111)_2=(\mathrm{D7})_{16}\)
(\cref{sec:fi-esadecimale-ottale}).

\begin{exercise}
\label{ex:fi-rappresentazione-ottale}
Che cos'è la rappresentazione ottale e come si collega alla base binaria?
\end{exercise}

La notazione \emph{ottale} impiega la base otto e le cifre da
\(0\) a \(7\). Poiché \(8=2^3\), una cifra ottale corrisponde a
un gruppo di tre bit. La conversione diretta dal binario avviene
raggruppando i bit a tre a tre, partendo dal separatore
(\cref{sec:fi-esadecimale-ottale}).

\begin{exercise}
\label{ex:fi-rappresentazione-caratteri}
Come vengono rappresentati i caratteri in un computer?
\end{exercise}

Una codifica associa valori numerici ai caratteri di un repertorio.
Un \emph{punto di codice} identifica un carattere, mentre il
\emph{formato di codifica} stabilisce come rappresentare il punto
di codice mediante byte. Lo stesso carattere può quindi occupare
sequenze di byte differenti secondo il formato
(\cref{sec:fi-caratteri}).

\begin{exercise}
\label{ex:fi-rappresentazione-ascii-unicode}
Che differenza c'è tra ASCII, ASCII esteso e Unicode?
\end{exercise}

\emph{ASCII} comprende \(128\) codici rappresentabili su sette bit.
\emph{ASCII esteso} non identifica un unico standard, ma diverse
codifiche a otto bit, come ISO-8859-1 e Windows-1252.
\emph{Unicode} definisce invece un repertorio di punti di codice
molto più ampio; formati come UTF-8, UTF-16 e UTF-32 ne specificano
la codifica in memoria (\cref{sec:fi-caratteri,tab:fi-ascii}).
